\chapter{OBJETIVO}

\section{Objetivo Geral}

Desenvolver uma solução de \textit{software} nativa em linguagem C para a implementação de um AAS em sistemas embarcados, permitindo avaliar a sua eficiência, viabilidade e adequação operacional em dispositivos com restrição de recursos computacionais.

\section{Objetivos Específicos}

Para alcançar o objetivo geral, foram definidos os seguintes objetivos específicos:
\begin{alineas}
    \item mapear o estado da arte e as soluções de AAS para sistemas embarcados, identificando as abordagens atuais, bem como suas limitações e lacunas tecnológicas;
    \item estabelecer os requisitos funcionais e não funcionais da biblioteca proposta, orientando-se pelos modelos padronizados do AAS, pelas demandas do protocolo OPC UA e pelas restrições do microcontrolador;
    \item implementar as funcionalidades centrais da biblioteca proposta, integrando a estrutura do AAS às rotinas de aquisição de dados e comunicação de rede;
    \item desenvolver um protótipo físico focado na instrumentação de grandezas elétricas para servir como um ambiente de validação prático (AAS Tipo 2);
    \item avaliar o desempenho, a estabilidade e o consumo de recursos da solução implementada, atestando sua eficiência em relação a implementações herdadas de ambientes computacionais mais robustos.
\end{alineas}
